\section{Julia-Mengen und Fatou-Mengen}
\label{julia:sec:julia}
\kopfrechts{Julia-Mengen und Fatou-Mengen}%
Julia-Mengen und Fatou-Mengen entstehen aus der Überlegung was geschieht, wenn man eine komplexe Funktion $f$ wiederholt auf ein $z \in \mathbb{C}$ anwendet.
Hieraus ergibt sich für jedes $z$ eine Folge komplexer Zahlen
\begin{equation}
    z \rightarrow f(z) \rightarrow f(f(z)) \rightarrow \ldots \:,
    \label{julia:julia_folge}
\end{equation}
bzw.
\begin{equation*}
    z_{n + 1} = f(z_n) \quad \text{mit}\: n \in \mathbb{N}_0\: \text{und einem Startwert}\: z_0 \in \overline{\mathbb{C}} = \mathbb{C} \cup \infty.
\end{equation*}

Je nach Wahl des Startwerts $z_0$ kann diese Folge zwei Verhalten zeigen:
\begin{enumerate}
    \item Eine kleine Änderung des Startwerts führt zur praktisch derselben Folge. Der Wert $z_0$ wird der \emph{Fatou-Menge} $F_f$ von $f$ zugeordnet.
\index{Fatou-Menge}%
    \item Eine kleine Änderung des Startwerts führt zu einer komplett anderen Folge. In diesem Fall wird $z_0$ der \emph{Julia-Menge} $J_f$ von $f$ zugeordnet.
\index{Julia-Menge}%
\end{enumerate}
Innerhalb dieses Abschnitts sollen die obigen, provisorischen Definitionen formalisiert werden.

Die Begriffe ``praktisch dieselbe Folge'' und ``komplett andere Folge'' beziehen sich darauf, gegen welche Werte die Folgen gehen, wenn $n \to \infty$.
Weiter ist mit ``eine kleine Änderung'' gemeint, dass ein $\varepsilon \in \mathbb{C}$ mit $|\varepsilon|$ nahe bei $0$ zum Startwert $z_0$ addiert wird.
Wenn es für einen Startwert ein $\varepsilon$ gibt, so dass man eine komplett andere Folge erhält, so gehört der Startwert zur Julia-Menge, auch wenn es andere $\varepsilon$ gibt, für die dies nicht der Fall ist.

\subsection{Periodische Folgen}
Das Verhalten von solchen Folgen ist besonders dann interessant, wenn diese periodisch sind, bzw. $z_n = z_0$ für ein $n > 0$ gilt.
In dem Fall wird $z_0$ als \emph{periodischer Punkt} bezeichnet.
\index{periodischer Punkt}%
Ist $n$ die kleinste natürliche Zahl mit dieser Eigenschaft, so heisst $n$ die \emph{Periode} des Zyklus.
\index{Periode}%
Falls dies bereits für $n = 1$ zutrifft, wenn also $f(z_0) = z_0$ gilt, so ist $z_0$, nicht nur ein periodischer Punkt, sondern auch ein \emph{Fixpunkt} von $f$.
\index{Fixpunkt}%
Aus diesen Definitionen geht hervor, dass ein periodischer Punkt von $f$ mit Periode $n$ gleichzeitig ein Fixpunkt der $n$-fach zusammengesetzten Funktion $f^n$ ist.

Um die Komplexität der Fragestellung zu reduzieren, werden zunächst nur Fixpunkte betrachtet.

\subsubsection{Anziehende und abstossende Fixpunkte}
Am Beispiel der Funktion 
\begin{equation*}
    f(z) = z^2,
\end{equation*}
lassen sich zwei verschiedene Verhaltensweisen von Fixpunkten beobachten.

Durch $n$-maliges Anwenden erhält man $f^n(z) = z^{2^n}$.
Weiter hat die Funktion $f$ drei Fixpunkte: $0, 1$ und $\infty$, denn für diese Punkte gilt $f(z_0) = z_0$.

Entwickelt man nun die Folge \eqref{julia:julia_folge} ausgehend von den Startwerten $0$ und $\infty$, so erhält man $f^n(0) = 0$ bzw. $f^n(\infty) \to \infty$.
Nun werden diese Startwerte leicht verändert, dies geschieht mittels einem $\varepsilon \in \mathbb{C}$, wobei $|\varepsilon|$ nahe bei $0$ ist:
\begin{itemize}
    \item Für $z_0 = 0$: $f^n(0 + \varepsilon) = (0 + \varepsilon)^{2^n} \to 0$ mit $|\varepsilon| > 0$
    \item Für $z_0 \to \infty$: $f^n(\infty + \varepsilon) = (\infty + \varepsilon)^{2^n} \to \infty$ mit $|\varepsilon| > 0$
\end{itemize}
Obwohl die Startwerte verändert wurden, erhält man schlussendlich die gleichen Resultate, wie wenn man die Startwerte nicht verändert hätte.
Anders ausgedrückt führen kleine Änderungen der Startwerte zu praktisch denselben Folgen.
Es scheint, dass diese beiden Startwerte Folgen in ihrer Umgebung anziehen, sprich dass es sich um \emph{anziehende} Fixpunkte handelt.
\index{anziehend}%
\index{Fixpunkt!anziehend}%

Für den Startwert $1$ erhält man $f^n(1) = 1$, durch leichtes Verändern jedoch:
\begin{itemize}
    \item $f^n(1 + \varepsilon) = (1 + \varepsilon)^{2^n} \to \infty$
    \item $f^n(1 - \varepsilon) = (1 - \varepsilon)^{2^n} \to 0$
\end{itemize}
Hier scheint es also, dass eine kleine Änderung des Startwerts zu einer ganz anderen Folge führt, heisst es handelt sich um einen \emph{abstossenden} Fixpunkt.
\index{Fixpunkt!abstossend}%

\subsubsection{Indifferente Fixpunkte}
Als zweites Beispiel dient die Funktion
\begin{equation*}
    h(z) = z + z^7,
\end{equation*}
bei der sich indifferente Fixpunkte beobachten lassen.
Diese Funktion hat Fixpunkte bei $-\infty, 0$ und $\infty$.
Die Fixpunkte $-\infty$ und $\infty$ verhalten sich erneut anziehend und werden daher nicht näher untersucht.
Nimmt man jedoch $0$ als Startwert, passiert etwas anderes.
Zunächst ergibt sich $h^n(0) = 0$, durch eine kleine Änderung erhält man,
\begin{equation*}
    h(0 + \varepsilon) = (0 + \varepsilon) + (0 + \varepsilon)^7.
\end{equation*}
Da es sich bei $|\varepsilon|$ um eine sehr kleine Zahl handelt, ist der zweite Summand gegenüber $\varepsilon$ vernachlässigbar, übrig bleibt der lineare Term.
Somit ist
\begin{equation*}
    h(0 + \varepsilon) \approx 0 + \varepsilon.
\end{equation*}
Wendet man nun erneut $h$ auf dieses Resultat an, erhält man wieder $0 + \varepsilon$.
Auch durch $n$-maliges Anwenden von $h$ ändert sich das Resultat nicht.
Die Folge ``bewegt'' sich also, anders als vorhin, nicht.
Der Startwert $0$, scheint also Folgen weder abzustossen, noch anzuziehen, man bezeichnet den Fixpunkt daher als \emph{indifferent}.
\index{Fixpunkt!indifferent}%

\subsubsection{Arten von Fixpunkten}
\label{julia:arten_periodische_punkte}
Das Verhalten einer Folge hängt also von der Art des verwendeten Fixpunktes ab.
Sei $z_0$ ein periodischer Punkt von $f$ mit Periode $n$, demzufolge ist $z_0$ ein Fixpunkt von $g = f^n$.
Dadurch wird die Frage nach dem Verhalten in der Nähe eines periodischen Punktes reduziert auf die Frage nach dem Verhalten in der Nähe eines Fixpunktes.
Diese Frage lässt sich mithilfe der Taylor-Reihe um $z_0$
\begin{equation*}
    g(z) \approx z_0 + g^\prime(z_0) (z - z_0)
    \label{julia:taylor}
\end{equation*}
beantworten, wobei höhere Potenzen in der Nähe von $z_0$ vernachlässigbar sind.
Verwendet man $z_0$ als Startwert und ändert diesen leicht, bedeutet dies
\begin{equation*}
    z = z_0 + \varepsilon.
\end{equation*}
Setzt man dies in die Taylor-Reihe ein, erhält man
\begin{equation*}
    g(z) = g(z_0 + \varepsilon) \approx z_0 + g^\prime(z_0) \left((z_0 + \varepsilon) - z_0\right) = z_0 + g^\prime(z_0) \varepsilon.
\end{equation*}
Wird $g$ erneut und wiederholt auf dieses Resultat angewendet, so wird das Verhalten der entstehenden Folge einzig von $g^\prime(z_0)$ festgelegt.
Ist $|g^\prime(z_0)| < 1$, so wird der zweite Summand durch wiederholte Anwendungen immer kleiner und die Folge kehrt zu $z_0$ zurück.
Bei $|g^\prime(z_0)| = 0$ geschieht dies bereits nach der ersten Anwendung von $g$.
Im Fall dass $|g^\prime(z_0)| > 1$ ist, wird der zweite Summand immer grösser und die Folge kehrt nicht zu $z_0$ zurück, sondern konvergiert zu einem anderen Wert.

Wenn $|g^\prime(z_0)| = 1$ ist, ergibt sich ein Spezialfall.
Durch wiederholtes Anwenden scheint das Resultat immer $z_0 + \varepsilon$ zu bleiben, das stimmt aber nicht.
Die Formel \eqref{julia:taylor} für die Taylor-Reihe ignoriert die höheren Potenzen, da diese in den drei anderen Fällen nichts am Ausgang ändern.
Wiederholtes Anwenden der Funktion führt dazu, dass die höheren Potenzen irgendwann nicht mehr vernachlässigt werden dürfen, und im Fall $|g^\prime(z_0)| = 1$ werden sie sogar sehr wichtig.
Das Verhalten der Folge wird hier durch diese höheren Potenzen bestimmt.
Die Folge kann sich so verhalten, wie eine Folge ausgehend von einem \emph{anziehenden} oder \emph{abstossenden} Fixpunkt, sie kann sich aber auch gar nicht ``bewegen''.
Dieses Phänomen wird im Abschnitt \ref{julia:sec:indifferent} näher beschrieben.

\begin{definition}
Der \emph{Multiplikator} eines Fixpunkts ist 
\index{Multiplikator}%
\begin{equation*}
    \lambda = g^\prime(z_0) = \left(f^n\right)^\prime(z_0).
    \label{julia:lambda_multiplicator}
\end{equation*}
Weiter sei $z_0$ ein periodischer Punkt,
dann heisst $z_0$ \emph{stark anziehend}, falls $\lambda = 0$,
\index{stark anziehend}%
\emph{anziehend}, falls $0 < |\lambda| < 1$,
\index{anziehend}%
\emph{indifferent}, falls $|\lambda| = 1$,
\index{indifferent}%
und \emph{abstossend}, falls $|\lambda| > 1$.
\index{abstossend}%
\label{julia:lambda_defs}
\end{definition}

Wendet man dies auf das vorherige Beispiel der Funktion $f(z) = z^2$ bzw. $f^n(z) = z^{2^n}$ an, bekommt man
\begin{equation*}
    \lambda = \left(f^n\right)^\prime(z_0) = 2^n z_0^{2^n - 1}.
\end{equation*}
Für die Fixpunkte von $f$ ergibt dies
\begin{itemize}
    \item $\lambda_0 = 0$,
    \item $\lambda_1 = 2^n$.
\end{itemize}
Somit sind diese Punkte stark anziehend, bzw.~abstossed.

Im Fall des Fixpunkts $\infty$ ist zunächst ein Koordinatenwechsel nötig, da sich dieser ausserhalb der komplexen Zahlenebene auf der Riemannschen Zahlenkugel befindet.
Sei
\begin{equation*}
    w = \frac{1}{z},
\end{equation*}
damit ist
\begin{equation*}
    f(w) = w^2 = \frac{1}{z^2}.
\end{equation*}
Für $w \to 0$ verhält sich $f$ gleich, wie für $z \to \infty$.
Berechnet man nun $\lambda$ unter Berücksichtigung der neuen Koordinate bekommt man
\begin{equation*}
    \lambda = \left(f^n\right)^\prime(w_0) = 2 w_0,
\end{equation*}
was für $w \to 0$, den Wert $0$ ergibt.
Insofern ist auch dieser Fixpunkt stark anziehend.

Die Namensgebungen decken sich also mit dem beobachteten Verhalten von Folgen, ausgehend von diesen Startwerten.
Es scheint, als ob eine Definition für die Julia-Mengen auf abstossenden Fixpunkten aufbauen muss. 

\subsection{Beobachtungen}
\label{julia:sec:beobachtungen}
Wie die vorhergehenden Beispiele zeigen, scheinen, gemäss den provisorischen Definition am Anfang des Abschnitts \ref{julia:sec:julia}, abstossende Fixpunkte zur Julia-Menge der jeweiligen Funktion zu gehören.
Ebenfalls scheinen anziehende Fixpunkte zur Fatou-Menge zu gehören.
Dies reicht jedoch noch nicht zur Definition dieser Mengen aus, denn indifferente Punkte, nicht-fixe periodische Punkte und nicht-periodische Punkte wurden noch nicht genauer betrachtet.

\subsubsection{Indifferente Punkte}
\label{julia:sec:indifferent}
Im Abschnitt \ref{julia:arten_periodische_punkte} wurde beschrieben, dass indifferente periodische Punkte einen Spezialfall bilden.
Abhängig von den höheren Potenzen in der Taylor-Reihe können indifferente Punkte entweder zur Julia-Menge oder zur Fatou-Menge ihrer jeweiligen Funktion gehören.
Wenn diese höheren Potenzen bei $n$-maliger Anwendung von $f$ verschwinden, gehört der indifferente Punkt zur Fatou-Menge.
Wachsen diese jedoch langsam, so verhält sich der indifferente Punkt wie ein abstossender Punkt und gehört daher zur Julia-Menge.
Weitere Details sind im Abschnitt 13 in \cite{julia:indifferente_punkte} beschrieben.

\subsubsection{Nicht-fixe periodische Punkte}
\label{julia:sec:nicht_fix_periodisch}
Bisher wurden nur Fixpunkte, also Punkte mit $f(z_0) = z_0$, und deren Verhalten untersucht.
Nicht-fixe periodische Punkte --- also Punkte mit $f^n(z_0) = z_0$, mit $n > 1$ --- verhalten sich grundsätzlich gleich, auch hier gibt es anziehende, abstossende und indifferente Punkte.

\begin{beispiel}
Nachfolgend wird dies am Beispiel der Funktion $f(z) = z^2$ illustriert.
Am Ende des Abschnitts \ref{julia:arten_periodische_punkte} wurde festgestellt, dass $f$ einen abstossenden Fixpunkt bei $1$ hat.
Diese Funktion hat jedoch weitere abstossende, nicht-fixe periodische Punkte --- wenn $n \to \infty$, sogar unendlich viele.

Damit ein Punkt periodisch, aber nicht fix ist, muss
\begin{equation*}
    f^n(z) = z^{2^n} = z,
\end{equation*}
für ein $n > 1$ gelten.
Betrachtet man die beiden rechten Ausdrücke dieser Gleichung und faktorisiert diese, erhält man
\begin{equation*}
    z \left(z^{2^n - 1} - 1\right) = 0.
\end{equation*}
Periodische Punkte von $f$ sind also:
\begin{itemize}
    \item $z = 0$, wie bereits beschrieben ist dieser Punkt stark anziehend und wird daher nicht weiter betrachtet.
    \item $z^{2^n - 1} = 1$.
\end{itemize}

Bei der zweiten Lösung handelt es sich um Einheitswurzeln $z_e$ des Grades $2^n - 1$, d. h. komplexe Zahlen, welche durch Potenzieren irgendwann $1$ ergeben.
Oder anders formuliert, um Wurzeln dieses Polynoms $2^n - 1$-ten Grades.
Für $n \to \infty$ gibt es unendlich viele Einheitswurzeln, was bedeutet, dass $f$ unendlich viele periodische Punkte hat.

Die Einheitswurzeln $z_e$ sind auch abstossend. Um dies zu zeigen berechnet man
\begin{equation*}
    \lambda = \left(f^n\right)^\prime(z_e) = 2^n z_e^{2^n - 1}.
\end{equation*}
Der zweite Faktor ist, aufgrund der Periodizität, gleich $1$, es bleibt
\begin{equation*}
    \lambda = 2^n 1 = 2^n
\end{equation*}
und weil $\lambda > 1$, handelt es sich um abstossende periodische Punkte.

Die Einheitswurzeln $z_e$ liegen stets auf dem Einheitskreis, d.~h. $|z_e| = 1$.
Wäre dies nicht der Fall, würde der Betrag durch das Potenzieren verschwinden, bzw. ansteigen und somit würde $z_e^{2^n - 1} = 1$ nicht mehr gelten.
Addiert man ein $\varepsilon$ mit Betrag nahe bei $0$ zu $z_e$, ist der Betrag entweder kleiner oder grösser als $1$.
Die Folge, die durch wiederholtes Anwenden von $f$ entsteht, wird nun entweder gegen $0$ oder gegen $\infty$ gehen.
Ohne $\varepsilon$ bewirkt $f$, dass das Argument sich um den Einheitskreis dreht und periodisch zu $1$ zurückkehrt.

Gemäss den provisorischen Definitionen am Anfang des Abschnitts \ref{julia:sec:julia} gehören die Einheitswurzeln daher in die Julia-Menge von $f$.
\end{beispiel}

Somit sind nicht nur abstossende Fixpunkte, sondern auch abstossende periodische Punkte relevant zur Definition der Julia-Mengen.

\subsubsection{Nicht-periodische Punkte}
\label{julia:sec:nicht_periodisch}
Punkte die nicht periodisch sind, können entweder zur Julia-Menge oder zur Fatou-Menge gehören.

\begin{beispiel}
Der Punkt $z_0 = e^{2\pi i\sqrt{2}}$ gehört zur Julia-Menge der Funktion $f(z) = z^2 = r e^{2\pi i(2 \varphi)}$, obwohl er nicht periodisch ist.
Durch $n$-maliges Anwenden wird $f$ zu $f^n(z) = z^{2^n} = r^{2^n} e^{2\pi i(2^n \varphi)}$.
Setzt man den Punkt $z_0$ ein, d.h. setzt $\sqrt{2}$ für $\varphi$ ein, müsste
\begin{equation*}
    2^n \sqrt{2} \equiv \sqrt{2} \mod{2\pi}
\end{equation*}
für ein $n$ gelten, damit $z_0$ periodisch ist.
Dies ist aber nur für $n = 0$ der Fall, was bedeutet dass $z_0$ nicht periodisch ist.

Nimmt man $z_0$ als Startwert für eine Folge und verändert diesen leicht, so gibt es drei Fälle
\begin{itemize}
    \item $|z_0 + \varepsilon| = 1$,
    \item $|z_0 + \varepsilon| < 1$,
    \item $|z_0 + \varepsilon| > 1$.
\end{itemize}
Wendet man nun wiederholt $f$ auf den veränderten Startwert an, so wird sich dieser im ersten Fall um den Einheitskreis ``drehen'' und ist möglicherweise sogar periodisch.
Im zweiten Fall wird der Betrag gegen $0$ gehen und im dritten Fall gegen $\infty$.
Je nach Wert von $\varepsilon$ erhält man unterschiedliche Folgen, das Verhalten ist chaotisch und somit gehört $z_0$ zur Julia-Menge von $f$.

Verwendet man für die gleiche Funktion jedoch $0.5$ als Startwert und verändert diesen leicht, so erhält man $f^n(0.5 + \varepsilon) = (0.5 + \varepsilon)^{2^n}$, was gegen $0$ konvergiert.
Die Folge für den unveränderten Startwert konvergiert ebenfalls gegen $0$, was bedeutet, dass dieser Punkt zur Fatou-Menge von $f$ gehört.
\label{julia:nicht_periodisch_julia}
\end{beispiel}

Nicht-periodische Punkte können sich also auch unterschiedlich verhalten und gehören daher entweder zur Julia-Menge oder zur Fatou-Menge.

\subsection{Definition der Julia-Mengen}
Wie im Abschnitt \ref{julia:sec:beobachtungen} gezeigt wurde, scheint es, dass eine Definition für Julia-Mengen
\begin{itemize}
    \item auf abstossenden periodischen Punkten (nicht nur Fixpunkten) aufbauen muss
    \item sowie indifferente periodische Punkte und nicht-periodische Punkte berücksichtigen muss.
\end{itemize}

\begin{beispiel}
Betrachtet man erneut die Funktion $f(z) = z^2$, so ist bereits bekannt, dass alle abstossenden Punkte in der Julia-Menge enthalten sind.
Alle diese Punkte liegen auf dem Einheitskreis und wie im Abschnitt \ref{julia:sec:nicht_fix_periodisch} beschrieben, gibt es unendlich viele davon.

Diese abstossenden Punkte sind jedoch nicht deckungsgleich mit dem Einheitskreis, es gibt ``Lücken'' auf dem Einheitskreis, welche nicht zu dieser Punktemenge gehören (siehe Beispiel \ref{julia:nicht_periodisch_julia}).
Allerdings ist die Distanz von jeder Lücke zu einem abstossenden Punkt beliebig klein.
Obwohl die ``Lückenpunkte'' selber nicht periodisch sind, verhalten sie sich gleich wie abstossende periodische Punkte weil sie so nahe bei diesen liegen.

Entwickelt man die Folge \eqref{julia:julia_folge} ausgehend von einem ``Lückenpunkt'' $z_l$ und entwickelt dann eine zweite Folge ausgehend von $z_l + \varepsilon$, erhält man ebenfalls zwei komplett andere Folgen.
Für den Startwert $z_l$ führt wiederholtes Anwenden von $f$ auch dazu, dass sich das Argument um den Einheitskreis dreht, jedoch nie mehr zu $1$ zurückkehrt.
Für $z_l + \varepsilon$ geht wiederholtes Anwenden gegen $0$, wenn $|\varepsilon|$ negativ ist und gegen $\infty$, wenn $|\varepsilon|$ positiv ist.
Diese ``Lückenpunkte'' gehören somit auch zur Julia-Menge.

Was die indifferenten periodischen Punkte angeht, so hat die Funktion $f$ keine davon.
Allgemein verhalten sich diese jedoch auch abstossend, wenn sie in einer solchen ``Lücke'' liegen.
Ihr genaues Verhalten ist aber wesentlich komplexer und zur vollständigen Beschreibung wird erneut auf Abschnitt 13 in \cite{julia:indifferente_punkte} verwiesen.

Abschliessend bedeutet das, dass die Julia-Menge von $f$ dem Rand des Einheitskreises entsprechen muss:
\begin{equation*}
    J_f = \left\{z \in \overline{\mathbb{C}}\:\mid\:|z| = 1\right\}. \qquad \qedhere
\end{equation*}
\end{beispiel}

\begin{definition}
Im Allgemeinen lässt sich die \emph{Julia-Menge} durch
\index{Julia-Menge}%
\begin{equation*}
    J_f \coloneqq \overline{\left\{z \in \overline{\mathbb{C}}\:\mid\:z\text{ ist ein abstossender periodischer Punkt von f}\right\}}
\end{equation*}
definieren.
Dabei bezeichnet $\overline{A}$ den topologischen Abschluss \cite{julia:abschluss} der Punktemenge $A$.
In diesem Fall bewirkt dies, dass die ``Lückenpunkte'' auch in $J_f$ enthalten sind.
\end{definition}

\subsection{Definition der Fatou-Mengen}
Die Definition der Fatou-Mengen ergibt sich unmittelbar aus der Definition der Julia-Mengen.
\begin{definition}
Da ein Punkt in der komplexen Ebene entweder zur Julia-Menge oder zur Fatou-Menge gehört, ist die \emph{Fatou-Menge} das Komplement der Julia-Menge
\index{Fatou-Menge}%
\begin{equation*}
    F_f \coloneqq \overline{\mathbb{C}} \setminus J_f.
\end{equation*}
\end{definition}

\begin{beispiel}
Im vorherigen Beispiel bestand die Julia-Menge aus dem Rand des Einheitskreises, folglich ist
\begin{equation*}
    F_f = \left\{z \in \overline{\mathbb{C}}\:\mid\:|z| \neq 1\right\}
\end{equation*}
die Fatou-Menge der Funktion $f(z) = z^2$.
\end{beispiel}

